\textbf{Phần 1. Trắc nghiệm lựa chọn}
\setcounterpageref{ex}{0}

\begin{ex} %[2X13N64]
Việc đếm số học sinh nam và số học sinh nữ trong lớp học được gọi là
\choice
{Vẽ tranh}
{\True Thu thập, phân loại và kiểm đếm số liệu}
{Giải toán có lời văn}
{Tính nhẩm}
\loigiai{
Đó là hoạt động thu thập, phân loại và kiểm đếm số liệu.\\
\textbf{Đáp án của câu hỏi là B}
}
\end{ex}

\begin{ex} %[2X13N64]
Trong một hộp có 5 viên bi đỏ, 4 viên bi xanh và 3 viên bi vàng. Các viên bi được phân loại theo
\choice
{Kích thước}
{Khối lượng}
{\True Màu sắc}
{Chất liệu}
\loigiai{
Phân loại theo màu sắc (đỏ, xanh, vàng).\\
\textbf{Đáp án của câu hỏi là C}
}
\end{ex}

\begin{ex} %[2X13H64]
Kiểm đếm các đồ dùng học tập của Lan: 4 cái bút chì, 2 cục tẩy, 1 cái thước kẻ. Lan có tất cả số đồ dùng là
\choice
{\True $7$ cái}
{$6$ cái}
{$8$ cái}
{$9$ cái}
\loigiai{
$4 + 2 + 1 = 7$ đồ dùng.\\
\textbf{Đáp án của câu hỏi là A}
}
\end{ex}

\begin{ex} %[2X13H64]
Kiểm đếm số cây trong vườn: 15 cây cam, 12 cây xoài, 8 cây bưởi. Loại cây nào có số lượng nhiều nhất?
\choice
{Cây xoài}
{Cây bưởi}
{Các loại bằng nhau}
{\True Cây cam ($15$ cây)}
\loigiai{
Cây cam có 15 cây là nhiều nhất.\\
\textbf{Đáp án của câu hỏi là D}
}
\end{ex}

\begin{ex} %[2X13H64]
Kiểm đếm số cây trong vườn: 15 cây cam, 12 cây xoài, 8 cây bưởi. Loại cây nào có số lượng ít nhất?
\choice
{\True Cây bưởi ($8$ cây)}
{Cây cam}
{Cây xoài}
{Không có}
\loigiai{
Cây bưởi có 8 cây là ít nhất.\\
\textbf{Đáp án của câu hỏi là A}
}
\end{ex}

\begin{ex} %[2X13H64]
Số cây cam nhiều hơn số cây bưởi là
\choice
{$5$ cây}
{$6$ cây}
{$8$ cây}
{\True $7$ cây}
\loigiai{
$15 - 8 = 7$ cây.\\
\textbf{Đáp án của câu hỏi là D}
}
\end{ex}

\begin{ex} %[2X13V64]
Tổ 1 có 4 bạn thích bơi, 3 bạn thích đá bóng, 5 bạn thích cầu lông. Tổng số bạn trong tổ 1 là
\choice
{\True $12$ bạn}
{$11$ bạn}
{$13$ bạn}
{$14$ bạn}
\loigiai{
$4 + 3 + 5 = 12$ bạn.\\
\textbf{Đáp án của câu hỏi là A}
}
\end{ex}

\begin{ex} %[2X13V64]
Trong một đĩa trái cây có: 6 quả chuối, 4 quả cam, 5 quả táo. Số quả chuối nhiều hơn số quả cam là
\choice
{$1$ quả}
{$3$ quả}
{\True $2$ quả}
{$4$ quả}
\loigiai{
$6 - 4 = 2$ quả.\\
\textbf{Đáp án của câu hỏi là C}
}
\end{ex}

\begin{ex} %[2X13V64]
Thống kê điểm 10 của lớp 2A trong tuần: Thứ Hai: 8 điểm 10; Thứ Ba: 10 điểm 10; Thứ Tư: 7 điểm 10; Thứ Năm: 9 điểm 10; Thứ Sáu: 11 điểm 10. Ngày có nhiều điểm 10 nhất là
\choice
{Thứ Ba}
{\True Thứ Sáu}
{Thứ Hai}
{Thứ Năm}
\loigiai{
Thứ Sáu có 11 điểm 10 (nhiều nhất).\\
\textbf{Đáp án của câu hỏi là B}
}
\end{ex}

\begin{ex} %[2X13V64]
Tổng số điểm 10 của lớp 2A trong cả tuần đó là
\choice
{$43$ điểm}
{$47$ điểm}
{$50$ điểm}
{\True $45$ điểm}
\loigiai{
$8 + 10 + 7 + 9 + 11 = 45$ điểm 10.\\
\textbf{Đáp án của câu hỏi là D}
}
\end{ex}

\begin{ex} %[2X13V64]
Một cửa hàng bán xe đạp ghi chép số xe bán trong tuần: Xe màu đỏ: 8 chiếc, xe màu xanh: 6 chiếc, xe màu đen: 10 chiếc. Số xe màu đen nhiều hơn xe màu xanh là
\choice
{$3$ chiếc}
{$5$ chiếc}
{\True $4$ chiếc}
{$2$ chiếc}
\loigiai{
$10 - 6 = 4$ chiếc.\\
\textbf{Đáp án của câu hỏi là C}
}
\end{ex}

\begin{ex} %[2X13C64]
Trong hộp có các viên bi màu xanh, đỏ, vàng. Biết số bi đỏ là 8 viên, số bi xanh gấp đôi bi đỏ, số bi vàng ít hơn bi xanh 4 viên. Số bi vàng có trong hộp là
\choice
{$10$ viên}
{\True $12$ viên}
{$14$ viên}
{$16$ viên}
\loigiai{
Bi xanh: $8 \times 2 = 16$ viên.\\
Bi vàng: $16 - 4 = 12$ viên.\\
\textbf{Đáp án của câu hỏi là B}
}
\end{ex}

\textbf{Phần 2. Trắc nghiệm đúng sai}
\setcounterpageref{ex}{0}

\begin{ex} %[2X13H64]
Một giá sách có: 18 quyển sách Toán, 15 quyển sách Tiếng Việt, 12 quyển truyện tranh. Xét tính đúng sai:
\choiceTF
{\True Sách Toán có số lượng nhiều nhất (18 quyển).}
{\True Truyện tranh có số lượng ít nhất (12 quyển).}
{Giá sách có tất cả 40 quyển.}
{\True Giá sách có tất cả 45 quyển ($18 + 15 + 12 = 45$).}
\loigiai{
\begin{itemize}
    \item Toán nhiều nhất: Đúng.
    \item Truyện ít nhất: Đúng.
    \item Tổng: $18 + 15 + 12 = 45 \neq 40$. Ý 3 Sai, Ý 4 Đúng.
\end{itemize}
\begin{itemize}
    \item a) \textbf{Đúng}.
    \item b) \textbf{Đúng}.
    \item c) \textbf{Sai}.
    \item d) \textbf{Đúng}.
\end{itemize}
}
\end{ex}

\begin{ex} %[2X13V64]
Kiểm đếm các phương tiện giao thông đi qua ngã tư trong 5 phút: 25 xe máy, 10 ô tô, 8 xe đạp. Xét tính đúng sai:
\choiceTF
{\True Xe máy đi qua nhiều nhất.}
{\True Xe đạp đi qua ít nhất.}
{Số ô tô nhiều hơn số xe máy.}
{\True Số xe máy nhiều hơn số ô tô là 15 chiếc ($25 - 10 = 15$).}
\loigiai{
\begin{itemize}
    \item Xe máy nhiều nhất: Đúng.
    \item Xe đạp ít nhất: Đúng.
    \item Xe máy (25) nhiều hơn ô tô (10). Ý 3 Sai, Ý 4 Đúng.
\end{itemize}
\begin{itemize}
    \item a) \textbf{Đúng}.
    \item b) \textbf{Đúng}.
    \item c) \textbf{Sai}.
    \item d) \textbf{Đúng}.
\end{itemize}
}
\end{ex}

\textbf{Phần 3. Trả lời ngắn}
\setcounterpageref{ex}{0}

\begin{ex} %[2X13H64]
Kiểm đếm số vật nuôi của nhà An: 12 con gà, 8 con vịt, 4 con thỏ. Tính tổng số vật nuôi nhà An.
\shortans{24}
\loigiai{
\textbf{Đáp số:} $24$\\
$12 + 8 + 4 = 24$ con.
}
\end{ex}

\begin{ex} %[2X13H64]
Số gà nhiều hơn số thỏ bao nhiêu con?
\shortans{8}
\loigiai{
\textbf{Đáp số:} $8$\\
$12 - 4 = 8$ con.
}
\end{ex}

\begin{ex} %[2X13V64]
Kiểm đếm số học sinh tham gia câu lạc bộ: Cờ vua có 8 bạn, Võ thuật có 12 bạn, Mĩ thuật có 10 bạn. Tính tổng số bạn tham gia 3 câu lạc bộ.
\shortans{30}
\loigiai{
\textbf{Đáp số:} $30$\\
$8 + 12 + 10 = 30$ bạn.
}
\end{ex}

\begin{ex} %[2X13V64]
Số bạn tham gia câu lạc bộ Võ thuật nhiều hơn Cờ vua bao nhiêu bạn?
\shortans{4}
\loigiai{
\textbf{Đáp số:} $4$\\
$12 - 8 = 4$ bạn.
}
\end{ex}

\begin{ex} %[2X13V64]
Trong tủ có 15 áo sơ mi trắng, 10 áo thun xanh và 5 áo khoác. Hỏi có tất cả bao nhiêu cái áo?
\shortans{30}
\loigiai{
\textbf{Đáp số:} $30$\\
$15 + 10 + 5 = 30$ cái áo.
}
\end{ex}

\begin{ex} %[2X13C64]
Trong một hộp bút có 3 loại bút: bút mực, bút chì và bút dạ. Biết có 6 bút mực, số bút chì gấp đôi bút mực, số bút dạ bằng một nửa số bút mực. Tính tổng số bút trong hộp.
\shortans{21}
\loigiai{
\textbf{Đáp số:} $21$\\
Bút chì: $6 \times 2 = 12$. Bút dạ: $6 : 2 = 3$.\\
Tổng: $6 + 12 + 3 = 21$ cái bút.
}
\end{ex}
